\hypertarget{appendix-5caa058bd9c7}{}\label{referencestart:appendix-5caa058bd9c7}
\section{Triple products and signed volume}\label{appendix-5caa058bd9c7-triple-products-and-signed-volume}

\hypertarget{appendix-5caa058bd9c7-triple-products-and-signed-volume}{}

\textbf{Curriculum depth:} Introduced · \textbf{First seen:} Phase6 · \textbf{Latest development:} Phase6

\begin{allmightroledefinition}

\subsection{Definition and mechanism}\label{appendix-5caa058bd9c7-definition-and-mechanism}

\hypertarget{appendix-5caa058bd9c7-definition-and-mechanism}{}

A scalar triple product measures oriented parallelepiped volume. Four-point orientation uses edge differences; tetrahedral volume is one sixth of the absolute result. A vector triple product is a different, vector-valued identity, and cross products are not associative.

\end{allmightroledefinition}

The panel with \(D=(0,0,2)\) has signed volume \(24\) and tetrahedron volume \(4\). Parenthesis changes can turn a nested cross from a nonzero vector into zero.

\subsection{Prerequisites}\label{appendix-5caa058bd9c7-prerequisites}

\hypertarget{appendix-5caa058bd9c7-prerequisites}{}

\begin{itemize}
\tightlist
\item
  \hyperref[appendix-2391db8052b1]{Cross products and oriented area} --- Oriented area supplies face normals, frame signs and geometric determinants.
\item
  \href{../06-GPU-Architecture-and-Real-Time-Rendering-P02/main.pdf\#nameddest=appendix-8bc4e8fc547c}{Dot products and orthogonality (GPU Architecture and Real-Time Rendering, Part 2, Reference notes: Dot products and orthogonality)} --- Inner-product measurements supply the perpendicularity and alignment constraints.
\end{itemize}

\subsection{Used by}\label{appendix-5caa058bd9c7-used-by}

\hypertarget{appendix-5caa058bd9c7-used-by}{}

\begin{itemize}
\tightlist
\item
  \hyperref[appendix-e1960925d1b8]{Exterior products and geometric algebra} --- Builds on this mechanism in the connected geometry story.
\end{itemize}

\subsection{Recall and next depth}\label{appendix-5caa058bd9c7-recall-and-next-depth}

\hypertarget{appendix-5caa058bd9c7-recall-and-next-depth}{}

Derive the vector identity by component expansion and preserve argument ordering. Full production proofs, target-specific integration and learner mastery require further evidence.

\subsection{Source and navigation}\label{appendix-5caa058bd9c7-source-and-navigation}

\hypertarget{appendix-5caa058bd9c7-source-and-navigation}{}

\href{../02-Mathematics-for-Visual-Computing-P01/main.pdf\#nameddest=01-foundations-i-ch017-signed-volume-extends-the-same-idea-into-space}{Teaching narrative (Mathematics for Visual Computing, Part 1, Ch. 1)} · \hyperref[appendix-fd090d982a15-phase-6---parts-501-600]{Coverage and source gap} · \hyperref[appendix-a103e7846d21-phase-6---source-analysis-and-verification]{Evidence and limits} · \hyperref[appendix-4ad47bdd269a]{Dependency Map}
